\scanpage{343}
\mspage{\unsure{9|1}}

\diagram{\begin{tikzcd}[ampersand replacement=\&,column sep=small,row sep=small]
{} \&[-.4em] {} \& {} \& b \& {} \\
\deleted{X}V={} \& \biggl(\unsure{0|o}\arrow[r,bend left=16,"\lambda"]\arrow[r,bend right=16,"\mu"'] \& a\arrow[ur,"\beta"]\arrow[dr,"\gamma"'] \& {} \& \biggr)\\
{} \& {} \& {} \& c
\end{tikzcd}}

avec
\[
\beta\lambda=\beta\mu\qquad\text{et}\qquad\gamma\lambda=\gamma\mu
\]
les deux seules flèches non identiques, en plus des \inserted{quatre} flèches marquées $\deleted{\unsure{B}}\beta,\deleted{\unsure{\beta}}\gamma,\deleted{\unsure{\beta}}\lambda,\mu$.

On trouve que \unsure{!|1}
\[
b\amalg_a c \overset{\sim}{\longrightarrow}e_X
\qquad\text{(somme amalgamée dans $\deleted{X}V^{\wedge}$)}
\]

Pourtant $\deleted{X}V$ n’est pas asphérique, ni même $1$-connexe : il \deleted{\unsure{a un|a le} \unsure{groupe|même}} \inserted{est} homotope \unsure{à|:} la sous-catégorie \unsure{ouverte|ouvert}
\diagram{\begin{tikzcd}[ampersand replacement=\&,column sep=small]
\deleted{X}V_0={} \& 0\arrow[r,bend left=18,"\lambda"]\arrow[r,bend right=18,"\mu"'] \& a,
\end{tikzcd}}
\deleted{\unsure{lorsq}} car on a une rétraction $\deleted{X}V\xrightarrow{r}\deleted{X}V_0$ (\unsure{envoie|envoyant} $b$ et $c$ sur $a$), et \deleted{$ir$} une flèche $ir\longrightarrow\mathrm{id}_{\deleted{X}V}$ (\unsure{et|où} $i:X_0\unsure{\hookrightarrow|\subset}X$ est l’inclusion). Donc $\deleted{X}V$ a le $=$ type d’homotopie que \uline{$\deleted{X}V_0$}, qui est un cercle.

Donc on a intérêt à prendre ce $V$, \inserted{\unsure{ou|où} $\widetilde V$ déduit de $V$ en rajoutant un élément final} \deleted{\unsure{comme}}
\diagram{\begin{tikzcd}[ampersand replacement=\&,column sep=small,row sep=small]
{} \& {} \& {} \& b\arrow[dr] \\
\widetilde V={} \& 0\arrow[r,bend left=18]\arrow[r,bend right=18] \& a\arrow[ur]\arrow[dr] \& {} \& e,\\
{} \& {} \& {} \& c\arrow[ur]
\end{tikzcd}}
comme $\unsure{W|W'}$, pour les contre-exemples page $6$ et $7$\deleted{\unsure{.}}

% archive sheet number: 343
